\section{Introduction}\label{sec:introduction}

The six-vertex model assigns an arrow to each edge of the square lattice,
subject to the rule that two arrows enter and two arrows leave each vertex.
Its height function is a discrete random surface. We prove the Gaussian
free field limit of this surface throughout the critical range
\(a=b=1\), \(0<c\le2\), for the plane state specified below.

\subsection{The plane law and the height convention}

The two vertices having both horizontal arrows directed inward, or both
directed outward, have weight \(c\); the other four vertices have weight
one. A finite configuration has the product of its vertex weights.
On
\[
 T_{M,L}=(\Z/M\Z)\times(\Z/L\Z),\qquad M,L\ge4,\quad L\ \text{even},
\]
condition every vertical column of horizontal edges to contain \(L/2\)
left arrows and \(L/2\) right arrows. This is a conserved sector of the row
transfer matrix. We first take \(M\to\infty\), and then even \(L\to\infty\),
in the local weak topology of arrow configurations. We denote the resulting
plane law by \(\P_c\). Existence in the previously untreated range is part
of the theorem.

Give the face on the left of an arrow height one more than the face on its
right. The ice rule makes this a well-defined height \(h\), modulo an
additive constant, on the faces of the plane. For \(x=(x_1,x_2)\in\R^2\),
let \(h(x)\) be the height of the face with bottom-left corner
\((\lfloor x_1\rfloor,\lfloor x_2\rfloor)\), and set
\[
 h_\delta(x)=h(x/\delta),\qquad \delta>0.
\]
This convention also fixes the value on grid lines.

Let \(\Gamma\) be the real plane Gaussian free field modulo constants,
normalized by the Green kernel
\begin{equation}\label{eq:green}
 G(x,y)=-\frac1{2\pi}\log|x-y|.
\end{equation}
Thus its covariance on smooth compactly supported tests of total integral
zero is the integral of \(G\) against their product.
For compactly supported finite signed measures \(\mu,\nu\) of total mass
zero, write
\[
 \mathcal E(\mu,\nu)=\iint G(x,y)\,\mu(\dd x)\nu(\dd y).
\]
We say that \(\mu\) has finite energy if
\[
 \iint \log_+\frac1{|x-y|}\,|\mu|(\dd x)|\mu|(\dd y)<\infty,
 \qquad \log_+t=\max(\log t,0).
\]
For compactly supported finite signed measures, this is equivalent to
requiring \(\mathcal E(\mu,\mu)\) to be an ordinary finite
signed-measure integral. In particular, the definition does not allow
cancellation between infinite positive and negative energies.

\begin{theorem}\label{thm:main}
For every fixed \(c\in(0,2]\), the prescribed iterated balanced-torus limit
\(\P_c\) exists. Define
\begin{equation}\label{eq:multiplier}
 \Delta=1-\frac{c^2}{2},\qquad
 \sigma(c)^2=\frac2{\arccos\Delta}=\frac1{\arcsin(c/2)}.
\end{equation}
As \(\delta\downarrow0\) through arbitrary positive real values,
the fields \(h_\delta\) converge to
\(\sigma(c)\Gamma\) in each of the following senses.
\begin{enumerate}
\item For every finite family \(\mu_1,\ldots,\mu_m\) of compactly supported
finite signed measures of total mass zero and finite energy,
\[
 \bigl(\langle h_\delta,\mu_j\rangle\bigr)_{j=1}^m
 \ \Longrightarrow\
 \mathcal N\!\left(0,
       \bigl(\sigma(c)^2\mathcal E(\mu_i,\mu_j)\bigr)_{i,j=1}^m\right).
\]
\item For every bounded open \(U\subset\R^2\) and \(-1<\alpha<0\),
there is convergence in law, with tightness, in the ordinary restriction
spaces
\[
 \mathcal C^\alpha_{\mathrm{res}}(U):=B^\alpha_{\infty,\infty}(U),
 \qquad B^\alpha_{p,q}(U),\qquad W^{\alpha,p}(U),
 \quad 1\le p<\infty,\quad 1\le q\le\infty,
\]
modulo constants. Convergence also holds in the outward completion
\eqref{eq:source-holder} defined below. No regularity of
\(\partial U\) is required. One may fix a smooth compactly supported
\(\rho\) of integral one and replace each field \(X\) by
\(X-\langle X,\rho\rangle\), using the same pinning for the limit.
\item For every \(k\ge1\) and every tuple
\((u_i,u_i')_{i=1}^k\) with
\(\{u_i,u_i'\}\cap\{u_j,u_j'\}=\varnothing\) for \(i\ne j\),
\begin{equation}\label{eq:increment-main}
 \E_c\prod_{i=1}^k
    \bigl(h_\delta(u_i')-h_\delta(u_i)\bigr)
 \longrightarrow
 \sigma(c)^k\sum_{\pi}\prod_{\{i,j\}\in\pi}C_{ij},
\end{equation}
where the sum is over pairings of \(\{1,\ldots,k\}\), is zero for odd \(k\),
and
\[
 C_{ij}=G(u_i',u_j')-G(u_i',u_j)-G(u_i,u_j')+G(u_i,u_j).
\]
The convergence is uniform on compact sets of such tuples.
\end{enumerate}
\end{theorem}

The theorem concerns this particular zero-slope plane state. It makes no
uniqueness assertion among all Gibbs states or all slopes, and no estimate
uniform as \(c\downarrow0\).

\paragraph{Regularity conventions.}
The restriction spaces in part~2 are restrictions of the corresponding
whole-plane distribution spaces, with their quotient norms. For
\(W^{\alpha,p}\) we include both the Bessel-potential and the
negative Slobodeckij realizations. The Hölder
laws are supported on the separable closure of smooth functions. This
choice makes the statement meaningful on every bounded open set without
an extension theorem for its boundary. We also prove convergence in the
intrinsic dual spaces \((W^{s,p'}(U))^*\), for \(0<s<1\) and
\(1\le p<\infty\), where the fractional seminorm uses
only pairs of points of \(U\); see Proposition~\ref{prop:intrinsic-dual}.

The outward convention is the literal seminorm of
\cite[Definition~2.6(iii)]{DKLM2026}. Let \(\mathcal T_1(Q)\), with
\(Q=(-1,1)^2\), consist of compactly supported signed densities in \(Q\),
of integral zero, whose zero extensions are \(1\)-Lipschitz. Set
\begin{equation}\label{eq:source-holder}
 \|f\|_{\mathrm{out},\alpha,U}
 =\sup_{\substack{0<\varepsilon\le1,\ \varphi\in\mathcal T_1(Q)\\
                   \supp\varphi\subset U/\varepsilon}}
   \varepsilon^{-\alpha}
   \left|\int f(x/\varepsilon)\varphi(x)\,\dd x\right|.
\end{equation}
Here \(U\) indexes the allowed tests of the global plane field;
the sampled points \(x/\varepsilon\) need not lie in \(U\).
The space is the completion of \(C_c^\infty(\R^2)\) modulo the
kernel of this seminorm.
The absolute value leaves the source convention unchanged because
the test family is symmetric. Its convergence requires a separate
argument from convergence in restriction spaces.

\subsection{Prior work and the new ingredients}

At \(c=1\), Lieb computed the exact square-ice entropy~\cite{Lieb1967}.
The Gaussian description of critical six-vertex heights and the prediction
of its coupling belong to the Coulomb-gas approach. In the F-model regime
\(1<c<2\), Di Francesco, Saleur and Zuber~\cite[Section~2.3]{DFSZ1987}
give a Gaussian-action prediction which, after converting their height
steps of \(\pi/2\) to unit steps, agrees with \eqref{eq:multiplier}.
This physics prediction is distinct from convergence in the measure
classes and distribution topologies of Theorem~\ref{thm:main}.

At the free-fermion point \(c=\sqrt2\), a dimer representation makes
the correlations determinantal. Kenyon's conformal-invariance and GFF
theorems for domino tilings~\cite{Kenyon2000,Kenyon2001} are foundational
scaling-limit results in this setting. The periodic dimer Gibbs measures,
surface tension and logarithmic fluctuations studied by Kenyon, Okounkov
and Sheffield~\cite{KOS2006} describe the corresponding plane-state
picture. Near the free-fermion point, Falco~\cite{Falco2013} derived
arrow-correlation asymptotics and logarithmic height variance for the
isotropic six-vertex model through an interacting-fermion representation.
Giuliani, Mastropietro and Toninelli proved GFF convergence for a symmetric
weakly interacting dimer model~\cite{GMT2017} and the relation between its
height amplitude and correlation exponent~\cite{GMTHaldane2017}.
Their later treatment of tilted profiles~\cite[Section~2.3 and Theorem~2]{GMT2020}
includes the six-vertex interaction and establishes height-covariance and
amplitude identities near the free-fermion point; its Remark~5 leaves the
higher-cumulant extension needed for a full GFF theorem undeveloped.

Duminil-Copin, Kozlowski, Krachun, Manolescu and
Tikhonovskaia~\cite{DCKKMT2022} proved existence and condensation of Bethe
roots and obtained the free energy and its expansion near half filling.
Duminil-Copin, Karrila, Manolescu and Oulamara~\cite{DCKMO2024} proved
logarithmic delocalization for \(1\le c\le2\), using crossing estimates
and free-energy information. Lis~\cite{Lis2021} developed the
Baxter--Kelland--Wu representation for multi-point spin correlations and
proved infinite-volume and delocalization results on subintervals of this
range. Glazman and Lammers~\cite{GlazmanLammers2025} subsequently obtained
delocalization, including logarithmic bounds in the isotropic case, by
percolation methods. These results establish the scale of fluctuations;
identifying their joint law requires additional information. The
variational framework of Sheffield~\cite{Sheffield2005} and the strict
convexity theorem of Lammers and Tassy~\cite{LammersTassy2024} provide the
equilibrium tools used below when \(c\ge1\).

Duminil-Copin, Kozlowski, Lammers and Manolescu~\cite{DKLM2026} give a
nonperturbative GFF theorem on \(\sqrt3\le c\le2\), with the exact
multiplier \eqref{eq:multiplier}. Their proof combines transfer-matrix
spectral representations and reflection positivity with rotational
invariance and regularity estimates. The rotational-invariance input is
listed as in preparation in the cited version. We retain the use of
transfer spectra and reflection, but for \(0<c<2\) construct the two
Gaussian currents from finite row identities and uniform thermodynamic
estimates. At \(c=2\), Appendix~\ref{sec:endpoint-rotation} supplies the
needed rotational-invariance statement from the independent fair
orientations of FK loops, the plane loop correspondence of~\cite{Lis2021},
and the homotopy coupling underlying the FK rotation theorem
of~\cite{DCKKMO2026Rotation}. The remaining endpoint
argument uses~\cite{DKLM2026}. Its conditional amplitude theorem is also
used on \(1\le c<2\), after convergence has been established. Thus the
full-range statement concerns precisely the balanced plane law defined
above, with its normalization fixed in physical height units.

The finite algebra is the algebraic Bethe ansatz. The norm determinant
goes back to Gaudin, McCoy and Wu~\cite{GaudinMcCoyWu1981} and
Korepin~\cite{Korepin1982}; the scalar-product formula is due to
Slavnov~\cite{Slavnov1989}. We use its trigonometric and diagonally
twisted forms as developed by Kitanine, Maillet and
Terras~\cite{KMT1999} and Kitanine, Maillet, Slavnov and
Terras~\cite{KMST2005}. These exact finite identities are the algebraic
input. Uniform control as the row length grows is a separate analytical
problem, addressed in Section~\ref{sec:thermodynamic}.

Three estimates connect this finite algebra to the limiting field.
Root-tail control and an edge maximum principle give uniform inverses
for Gaudin matrices and for mixed discrete--continuous quadrature
operators. Finite-dimensional residue conditions then control simultaneous
escapes of several rapidities in the determinant formulas. Next, bounds
on the zeros of angular spectral functions outside two separated scale
windows justify the limiting holomorphic and antiholomorphic currents.
Finally, when \(c<1\), a homogeneous stiffness lower bound and a
staggered vacuum-overlap upper bound identify the same covariance
coefficient. This last step supplies the normalization where the
conditional probabilistic input is unavailable. Moment estimates and
desmoothing pass from the currents to the test classes and separated
increments in Theorem~\ref{thm:main}.

\subsection{Proof overview}\label{sec:overview}

For \(0<c<2\), introduce
\[
 \lambda=\arccos(-\Delta)\in(0,\pi),
 \qquad \eta=\lambda/\pi.
\]
Section~\ref{sec:finite} defines commuting row operators and their diagonal
marks, which record height increments. Section~\ref{sec:thermodynamic}
constructs their limiting matrix elements, proves uniform tail bounds,
and shows that words with widely separated angular parameters decorrelate.
A positive full-line resolvent gives uniform density comparisons for
every fixed positive \(\lambda\), including the small-\(\lambda\) range.

These matrix elements define a Hilbert space with commuting angular rows
in Section~\ref{sec:angular}. Their joint spectral values are inner functions
of a complex row parameter. The finite folded-row continuation argument
in Lemma~\ref{lem:scalar-continuation} gives a common analytic neighborhood;
canonical inner-function factorization describes the possible zeros and
boundary singularities~\cite{AiDongJiang2019}. A mark acts as a current.
The finite mass
\(\kappa\) of the spectrum transverse to common angular dilation is the
coefficient of the logarithmic height covariance. Section~\ref{sec:normalization}
computes \(\kappa\) directly when \(c<1\); for \(c\ge1\) the cited conditional
normalization theorem applies once convergence has been established.

Section~\ref{sec:chiral} separates the angular spectrum into two scales and
proves the Wick rules for the limiting currents. Section~\ref{sec:probability}
passes from these currents to smooth tests, separated point increments
and finite-energy measures, then proves convergence in negative
regularity spaces. Its wavelet
argument uses compactly supported wavelets~\cite{Daubechies1988} and
the classical Besov coefficient characterization in the form of
\cite[Proposition~2.3]{HairerLabbe2017}; restriction from the plane avoids
any regularity assumption on the boundary of \(U\).
The final proof combines the interval \(0<c<2\) with the endpoint
argument of Appendix~\ref{sec:endpoint-rotation}. Each thermodynamic
passage specifies which size tends to infinity first.
